\documentclass[12pt,a4paper]{article}
\usepackage[utf8]{inputenc}
\usepackage[spanish,es-noquoting]{babel}
\usepackage[T1]{fontenc}
\usepackage{lmodern}
\usepackage{geometry}
\geometry{top=2cm,bottom=2cm,left=2.5cm,right=2.5cm}
\usepackage{listings}
\usepackage{fancyhdr}
\usepackage{graphicx}
\usepackage{tikz}
\usetikzlibrary{positioning}
\usepackage[most]{tcolorbox}
\newtcolorbox{ejerciciografico}{colback=white,colframe=black!65,boxrule=0.7pt,arc=2pt,left=8pt,right=8pt,top=7pt,bottom=7pt}
\graphicspath{{../../assets/}}
\lstset{language=C,basicstyle=\ttfamily\small,breaklines=true,frame=single,showstringspaces=false}
\setlength{\headheight}{15pt}
\pagestyle{fancy}
\fancyhf{}
\lhead{Monitorías Sistemas Operativos 2026-II}
\rhead{Capítulo 1: Procesos, nivel avanzado}
\cfoot{\thepage}
\title{\textbf{Taller de Lógica II: Bifurcaciones y Recursión}}
\author{Monitorías Sistemas Operativos}
\date{\today}
\begin{document}
\maketitle
\vspace{1.5cm}
\section*{Bloque 1: Diseño de jerarquías}
\begin{ejerciciografico}
\subsection*{Ejercicio 1}
\begin{center}\begin{tikzpicture}[node distance=.8cm,every node/.style={draw,circle,inner sep=2.2pt,fill=black}]
\node (r) {}; \node (a) [below left=of r] {}; \node (b) [below=of r] {}; \node (c) [below right=of r] {}; \draw (r)--(a) (r)--(b) (r)--(c);
\foreach \i in {1,...,3} {\node (n\i) [below=of b,xshift=\i*.8cm-1.6cm] {}; \draw (b)--(n\i);}
\end{tikzpicture}\end{center}
\end{ejerciciografico}
\begin{ejerciciografico}
\subsection*{Ejercicio 2}
\begin{center}\begin{tikzpicture}[node distance=.75cm,every node/.style={draw,circle,inner sep=2.2pt,fill=black}]
\node (r) {}; \node (a) [below left=of r] {}; \node (b) [below right=of r] {}; \draw (r)--(a) (r)--(b);
\node (a1) [below=of a] {}; \node (a2) [below=of a1] {}; \node (a3) [below=of a2] {}; \draw (a)--(a1)--(a2)--(a3);
\node (b1) [below left=of b] {}; \node (b2) [below right=of b] {}; \draw (b)--(b1) (b)--(b2);
\end{tikzpicture}\end{center}
\end{ejerciciografico}
\begin{ejerciciografico}
\subsection*{Ejercicio 3}
\begin{center}\begin{tikzpicture}[node distance=.7cm,every node/.style={draw,circle,inner sep=2.2pt,fill=black}]
\node (r) {}; \node (a) [below left=of r] {}; \node (b) [below right=of r] {}; \draw (r)--(a) (r)--(b);
\node (a1) [below left=of a] {}; \node (a2) [below right=of a] {};
\node (b1) [below left=of b] {}; \node (b2) [below right=of b] {};
\draw (a)--(a1) (a)--(a2) (b)--(b1) (b)--(b2);
\end{tikzpicture}\end{center}
\end{ejerciciografico}
\begin{ejerciciografico}
\subsection*{Ejercicio 4}
\begin{center}\begin{tikzpicture}[node distance=.8cm,every node/.style={draw,circle,inner sep=2.2pt,fill=black}]
\node (r) {}; \foreach \i in {0,...,4} {\node (h\i) [below=of r,xshift=\i*1.2cm-2.4cm] {}; \draw (r)--(h\i);}
\node (n0) [below=of h0] {}; \node (n3) [below=of h3] {}; \draw (h0)--(n0) (h3)--(n3);
\end{tikzpicture}\end{center}
\end{ejerciciografico}
\begin{ejerciciografico}
\subsection*{Ejercicio 5}
\begin{center}\begin{tikzpicture}[node distance=.8cm,every node/.style={draw,circle,inner sep=2.2pt,fill=black}]
\node (r) {}; \node (a) [below=of r] {}; \node (b) [below=of a] {}; \draw (r)--(a)--(b);
\foreach \i in {1,...,3} {\node (n\i) [below=of b,xshift=\i*.9cm-1.8cm] {}; \draw (b)--(n\i);}
\end{tikzpicture}\end{center}
\end{ejerciciografico}
\begin{ejerciciografico}
\subsection*{Ejercicio 6}
\begin{center}\begin{tikzpicture}[node distance=.75cm,every node/.style={draw,circle,inner sep=2.2pt,fill=black}]
\node (r) {}; \foreach \i in {1,...,4} {\node (h\i) [below=of r,xshift=\i*1.6cm-4cm] {}; \draw (r)--(h\i);}
\node (a1) [below=of h1] {}; \node (b1) [below=of a1] {}; \node (a4) [below=of h4] {}; \node (b4) [below=of a4] {}; \draw (h1)--(a1)--(b1) (h4)--(a4)--(b4);
\foreach \x in {2,3} {\node (l\x) [below left=of h\x] {}; \node (q\x) [below right=of h\x] {}; \draw (h\x)--(l\x) (h\x)--(q\x);}
\end{tikzpicture}\end{center}
\end{ejerciciografico}
\begin{ejerciciografico}
\subsection*{Ejercicio 7}
\begin{center}\begin{tikzpicture}[node distance=.7cm,every node/.style={draw,circle,inner sep=2.2pt,fill=black}]
\node (r) {}; \node (a) [below left=of r] {}; \node (x1) [below right=of r] {}; \draw (r)--(a) (r)--(x1);
\node (b) [below left=of a] {}; \node (x2) [below right=of a] {}; \draw (a)--(b) (a)--(x2);
\node (c) [below left=of b] {}; \node (x3) [below right=of b] {}; \draw (b)--(c) (b)--(x3);
\node (d) [below left=of c] {}; \node (x4) [below right=of c] {}; \draw (c)--(d) (c)--(x4);
\end{tikzpicture}\end{center}
\end{ejerciciografico}
\begin{ejerciciografico}
\subsection*{Ejercicio 8}
\begin{center}\begin{tikzpicture}[node distance=.75cm,every node/.style={draw,circle,inner sep=2.2pt,fill=black}]
\node (r) {}; \node (a) [below left=of r] {}; \node (b) [below right=of r] {}; \draw (r)--(a) (r)--(b);
\node (a1) [below=of a] {}; \node (b1) [below=of b] {}; \draw (a)--(a1) (b)--(b1);
\node (x) [below left=of b1] {}; \node (y) [below right=of b1] {}; \draw (b1)--(x) (b1)--(y);
\end{tikzpicture}\end{center}
\end{ejerciciografico}
\begin{ejerciciografico}
\subsection*{Ejercicio 9}
\begin{center}\begin{tikzpicture}[node distance=.75cm,every node/.style={draw,circle,inner sep=2.2pt,fill=black}]
\node (r) {}; \node (a) [below=of r] {}; \draw (r)--(a);
\node (r1) [below left=of r] {}; \node (a1) [below=of a] {}; \draw (r)--(r1) (a)--(a1);
\node (z) [below=of r1] {}; \draw (r1)--(z);
\end{tikzpicture}\end{center}
\end{ejerciciografico}
\begin{ejerciciografico}
\subsection*{Ejercicio 10}
\begin{center}\begin{tikzpicture}[node distance=.75cm,every node/.style={draw,circle,inner sep=2.2pt,fill=black}]
\node (r) {}; \foreach \i in {1,...,4} {\node (h\i) [below=of r,xshift=\i*1.4cm-3.5cm] {}; \draw (r)--(h\i);}
\foreach \i in {1,3} {\node (a\i) [below left=of h\i] {}; \node (b\i) [below right=of h\i] {}; \draw (h\i)--(a\i) (h\i)--(b\i);}
\end{tikzpicture}\end{center}
\end{ejerciciografico}
\section*{Bloque 2: Análisis de código}
\textit{Para cada caso, determine procesos totales, árbol y procesos que alcanzan el último \texttt{printf}.}
\subsection*{Ejercicio 11}
\begin{lstlisting}
for (int i = 0; i < 3; i++)
    if (fork() == 0) break;
printf("X\n");
\end{lstlisting}
\subsection*{Ejercicio 12}
\begin{lstlisting}
if (fork() != 0) {
    fork();
} else if (fork() == 0) {
    fork();
}
\end{lstlisting}
\subsection*{Ejercicio 13}
\begin{lstlisting}
pid_t p = fork();
if (p == 0) {
    for (int i = 0; i < 2; i++) fork();
}
\end{lstlisting}
\subsection*{Ejercicio 14}
\begin{lstlisting}
if (fork() || fork()) {
    if (fork() == 0) printf("A\n");
}
\end{lstlisting}
\subsection*{Ejercicio 15}
\begin{lstlisting}
for (int i = 0; i < 2; i++) {
    if (fork() != 0) break;
}
\end{lstlisting}
\subsection*{Ejercicio 16}
\begin{lstlisting}
switch (fork()) {
case 0: fork(); break;
default: if (fork() == 0) fork();
}
\end{lstlisting}
\subsection*{Ejercicio 17}
\begin{lstlisting}
pid_t a = fork();
if (a == 0 && fork() != 0) fork();
\end{lstlisting}
\subsection*{Ejercicio 18}
\begin{lstlisting}
if (fork() == 0) {
    fork();
} else {
    for (int i = 0; i < 2; i++) fork();
}
\end{lstlisting}
\subsection*{Ejercicio 19}
\begin{lstlisting}
for (int i = 0; i < 2; i++) {
    pid_t p = fork();
    if (p == 0) break;
}
\end{lstlisting}
\subsection*{Ejercicio 20}
\begin{lstlisting}
void crear(int n) {
    if (n <= 0) return;
    if (fork() == 0) crear(n - 1);
    else if (n == 2) fork();
}
int main(void) { crear(3); }
\end{lstlisting}
\end{document}
